% Exam answer booklets must only have the exam and ans options enabled
\documentclass[exam,ans]{yqteach}

%% Add in any additional settings needed here

%%%%%%%%%%%%%%%%%%%%%%%%%%%%%%%%%%%%%%%%%%%%%

% set the appropriate details for the module here
% Set details about the module here
\setinstitution{National University of Singapore}
\setdepartment{Department of Computer Science, School of Computing}
\setmodulecode{CS9999}
\setmoduletitle{Teaching Document Templates with \LaTeX}
\setfullacademicyear{2000/2001}
\setshortacademicyear{00/01}
\setsemester{Semester 2}
\setcopyrightyear{2023}

% Write the details of the exam here
\setdocumenttitle{Final Assessment}
\setduedate{1 January 2000}
\settimeallowed{1 hour}
\settotalscore{5}
%%%%%%%%%%%%%%%%%%%%%%%%%%%%%%%%%%%%

\begin{document}

\maketitle

% Use the markingsummary environment to generate the marking summary at the front page.
\begin{markingsummary}
    % Use the \markingsection command to enter a row in the marking summary table.
    % The first argument is the section name, the second number is the marks for that section.
    % Remember to set the total score (above) correctly.
    \markingsection{Multiple Choice Questions}{2}
    \markingsection{Short Answer Questions}{3}
\end{markingsummary}

\pagebreak

% Generate MCQ bubbles using the provided Python script and paste it in blank_mcq_tables.tex
\begin{table}[h]
\centering
\large
\bgroup
\def\arraystretch{1.5}
\begin{tabular}{|ccccc|ccccc|ccccc|}
\hline
&\textbf{A}&\textbf{B}&\textbf{C}&\textbf{D}&&\textbf{A}&\textbf{B}&\textbf{C}&\textbf{D}&&\textbf{A}&\textbf{B}&\textbf{C}&\textbf{D}\\
\hline
\textbf{1}&\begin{tikzpicture}[baseline=-3.5]
  \draw (0,0) circle (0.22cm);
\end{tikzpicture}&\begin{tikzpicture}[baseline=-3.5]
  \draw (0,0) circle (0.22cm);
\end{tikzpicture}&\begin{tikzpicture}[baseline=-3.5]
  \draw (0,0) circle (0.22cm);
\end{tikzpicture}&\begin{tikzpicture}[baseline=-3.5]
  \draw (0,0) circle (0.22cm);
\end{tikzpicture}&\textbf{2}&\begin{tikzpicture}[baseline=-3.5]
  \draw (0,0) circle (0.22cm);
\end{tikzpicture}&\begin{tikzpicture}[baseline=-3.5]
  \draw (0,0) circle (0.22cm);
\end{tikzpicture}&\begin{tikzpicture}[baseline=-3.5]
  \draw (0,0) circle (0.22cm);
\end{tikzpicture}&\begin{tikzpicture}[baseline=-3.5]
  \draw (0,0) circle (0.22cm);
\end{tikzpicture}&\textbf{3}&\begin{tikzpicture}[baseline=-3.5]
  \draw (0,0) circle (0.22cm);
\end{tikzpicture}&\begin{tikzpicture}[baseline=-3.5]
  \draw (0,0) circle (0.22cm);
\end{tikzpicture}&\begin{tikzpicture}[baseline=-3.5]
  \draw (0,0) circle (0.22cm);
\end{tikzpicture}&\begin{tikzpicture}[baseline=-3.5]
  \draw (0,0) circle (0.22cm);
\end{tikzpicture}\\
\hline
\textbf{4}&\begin{tikzpicture}[baseline=-3.5]
  \draw (0,0) circle (0.22cm);
\end{tikzpicture}&\begin{tikzpicture}[baseline=-3.5]
  \draw (0,0) circle (0.22cm);
\end{tikzpicture}&\begin{tikzpicture}[baseline=-3.5]
  \draw (0,0) circle (0.22cm);
\end{tikzpicture}&\begin{tikzpicture}[baseline=-3.5]
  \draw (0,0) circle (0.22cm);
\end{tikzpicture}&\textbf{5}&\begin{tikzpicture}[baseline=-3.5]
  \draw (0,0) circle (0.22cm);
\end{tikzpicture}&\begin{tikzpicture}[baseline=-3.5]
  \draw (0,0) circle (0.22cm);
\end{tikzpicture}&\begin{tikzpicture}[baseline=-3.5]
  \draw (0,0) circle (0.22cm);
\end{tikzpicture}&\begin{tikzpicture}[baseline=-3.5]
  \draw (0,0) circle (0.22cm);
\end{tikzpicture}&\textbf{6}&\begin{tikzpicture}[baseline=-3.5]
  \draw (0,0) circle (0.22cm);
\end{tikzpicture}&\begin{tikzpicture}[baseline=-3.5]
  \draw (0,0) circle (0.22cm);
\end{tikzpicture}&\begin{tikzpicture}[baseline=-3.5]
  \draw (0,0) circle (0.22cm);
\end{tikzpicture}&\begin{tikzpicture}[baseline=-3.5]
  \draw (0,0) circle (0.22cm);
\end{tikzpicture}\\
\hline
\textbf{7}&\begin{tikzpicture}[baseline=-3.5]
  \draw (0,0) circle (0.22cm);
\end{tikzpicture}&\begin{tikzpicture}[baseline=-3.5]
  \draw (0,0) circle (0.22cm);
\end{tikzpicture}&\begin{tikzpicture}[baseline=-3.5]
  \draw (0,0) circle (0.22cm);
\end{tikzpicture}&\begin{tikzpicture}[baseline=-3.5]
  \draw (0,0) circle (0.22cm);
\end{tikzpicture}&\textbf{8}&\begin{tikzpicture}[baseline=-3.5]
  \draw (0,0) circle (0.22cm);
\end{tikzpicture}&\begin{tikzpicture}[baseline=-3.5]
  \draw (0,0) circle (0.22cm);
\end{tikzpicture}&\begin{tikzpicture}[baseline=-3.5]
  \draw (0,0) circle (0.22cm);
\end{tikzpicture}&\begin{tikzpicture}[baseline=-3.5]
  \draw (0,0) circle (0.22cm);
\end{tikzpicture}&\textbf{9}&\begin{tikzpicture}[baseline=-3.5]
  \draw (0,0) circle (0.22cm);
\end{tikzpicture}&\begin{tikzpicture}[baseline=-3.5]
  \draw (0,0) circle (0.22cm);
\end{tikzpicture}&\begin{tikzpicture}[baseline=-3.5]
  \draw (0,0) circle (0.22cm);
\end{tikzpicture}&\begin{tikzpicture}[baseline=-3.5]
  \draw (0,0) circle (0.22cm);
\end{tikzpicture}\\
\hline
\textbf{10}&\begin{tikzpicture}[baseline=-3.5]
  \draw (0,0) circle (0.22cm);
\end{tikzpicture}&\begin{tikzpicture}[baseline=-3.5]
  \draw (0,0) circle (0.22cm);
\end{tikzpicture}&\begin{tikzpicture}[baseline=-3.5]
  \draw (0,0) circle (0.22cm);
\end{tikzpicture}&\begin{tikzpicture}[baseline=-3.5]
  \draw (0,0) circle (0.22cm);
\end{tikzpicture}&&&&&&&&&&\\
\hline
\end{tabular}
\egroup
\end{table}

% Use the \cq command to start a question while setting the question counter. For example, \cq[123]{456} resets the question counter to 123 and creates Question 123 [456 marks].
% The \csq command does the same except for sub-questions. For example, \csq[123][3]{456} creates Question 123 (iii) [456 marks].
\csq[11][1]{2} Tell Charlie he is wrong by giving him a scenario in which attempting to perform a well-formed SQL query to a database may still yield an exception.

% use \answerbox to create a simple box for students to write answers. The argument is the height of the box
\answerbox{500}

% use \extraspace to give the declaration that the following space is extra space.
\extraspace
\answerbox{620}

%%% Don't touch the following

\endofdocument

\end{document}
